\subsection{Construction of the Causal Lower Bound}
\label{sec:causal-bound-components}

The controller derives a positive idle-time certificate from three sources:
the instruction prefix before the next SA operation, the position of required
memory requests in the core request FIFO, and the number of required memory
responses that have not yet returned. Each quantity is converted into a lower
bound under maximum-rate progress. The controller never substitutes an
average memory latency or a subsequently observed completion time.

\subsubsection{Instruction-Prefix Bound}

For each resident tile, the simulator locates the first remaining
\texttt{GEMM\_PRELOAD} or \texttt{GEMM}. Let $P_i$ denote the number of
instructions preceding that candidate in tile $i$. Because the core can
successfully remove at most one resident instruction per cycle, the candidate
cannot be issued for at least
\begin{equation}
  B^{\mathrm{prefix}}_i=P_i
  \label{eq:prefix-bound}
\end{equation}
cycles. By construction, every instruction counted in $P_i$ is non-SA: the
scan stops at the first SA instruction. Therefore, the prefix itself cannot
require the powered-down array.

When $P_i>0$, the implemented controller uses the prefix directly and does
not add a memory bound. Operand state can change while the prefix drains, but
such changes cannot make the SA candidate pass the preceding instructions.
If the operands remain unavailable, execution starts later; if they become
ready early, the prefix remains binding. This choice preserves a valid lower
bound without assuming independence between instruction issue and memory
progress.

\subsubsection{Core-FIFO Ticket Bound}

A \texttt{MOVIN} is divided into one 32-byte \texttt{MemoryAccess} object per
source address. ONNXim assigns these objects monotonically increasing tickets
when they enter the core request FIFO. For each scratchpad allocation, the
predictor stores the ticket of the final request required by that allocation.
The core also records the ticket of the most recently dequeued request and
checks that dequeues preserve consecutive FIFO order.

Let $q_{i,j}$ be the tail ticket for source operand $j$ of tile $i$, and let
$d$ be the most recently dequeued ticket. When $q_{i,j}>d$, its current FIFO
position is
\begin{equation}
  Q_{i,j}=q_{i,j}-d.
  \label{eq:fifo-position}
\end{equation}
The evaluated configuration has 16 injection ports per core and equal 1-GHz
core and NoC clocks. Even under maximum service, the tail request cannot leave
the FIFO earlier than
\begin{equation}
  B^{\mathrm{fifo}}_{i,j}
    =\left\lfloor\frac{Q_{i,j}-1}{16}\right\rfloor
  \label{eq:fifo-ticket-bound}
\end{equation}
core cycles. For example, positions 1--16 produce a zero-cycle bound, while
position 17 produces a one-cycle bound. This equation is enabled only when
ticket order is exact, the operand is allocated but invalid, memory requests
remain pending, and the tail ticket is still ahead of the dequeue frontier.
Otherwise, this source contributes no positive certificate. The synthesized
implementation uses 20-bit ticket and horizon fields; values outside the
represented forward window are conservatively treated as unsupported rather
than extrapolated.

Equation~\ref{eq:fifo-ticket-bound} is configuration specific. It is exact
for the evaluated 16-port, equal-clock configuration because division by 16
is implemented as a four-bit right shift. A configuration with a different
port count or clock ratio requires a separately derived constant-ratio bound;
the evaluated RTL disables the specialized ticket certificate when its
configuration check is not satisfied.

\subsubsection{Outstanding-Response Bound}

Leaving the core FIFO does not make an operand ready. A read remains
outstanding while it traverses the NoC, DRAM, and return path. Each allocated
scratchpad entry therefore maintains a pending-memory count that is
initialized to the number of required requests and decremented only when a
read response fills that entry. Let $N_{i,j}$ be the sampled number of
required responses that have not yet returned for source operand $j$.

At most 16 responses for a core can become visible between consecutive
controller evaluations in the evaluated equal-clock configuration. Hence,
even under maximum response progress, the operand cannot become valid for at
least
\begin{equation}
  B^{\mathrm{resp}}_{i,j}
    =\left\lfloor\frac{N_{i,j}-1}{16}\right\rfloor
  \label{eq:response-drain-bound}
\end{equation}
cycles. The bound is emitted only for a used, allocated, invalid operand with
$N_{i,j}>0$. Counts are represented in 20 bits; saturation can only shorten
the certificate and is therefore conservative.

The pending count includes requests that are still in the core FIFO as well
as requests already traveling through the memory system. Consequently, the
FIFO and response-drain intervals may overlap. They are two lower bounds on
the same operand-ready time, not two independent serial delays. Adding them
could count the same requests twice and overestimate the idle interval. The
implemented controller therefore combines them using
\begin{equation}
  B_{i,j}=\max\!\left(B^{\mathrm{fifo}}_{i,j},
                       B^{\mathrm{resp}}_{i,j}\right).
  \label{eq:memory-bound-max}
\end{equation}
For a head SA candidate with two represented source operands, both operands
must become ready; therefore,
\begin{equation}
  B^{\mathrm{head}}_i=\max_j B_{i,j}.
  \label{eq:head-operand-bound}
\end{equation}
If an operand is already valid, it contributes zero. An unallocated or
otherwise unclassifiable operand does not receive a guessed positive memory
bound.

\subsubsection{Hardware-State Boundary}

The predictor-support state is not free. ONNXim adds a monotonically
increasing ticket to each generated memory request, a last-dequeued ticket per
core, a tail ticket per tracked scratchpad allocation, and pending-memory
counts used by the controller descriptor. The synthesized controller includes
the registered descriptor fields and the arithmetic that evaluates
Eqs.~\ref{eq:fifo-ticket-bound}--\ref{eq:head-operand-bound}. However, the
reported controller synthesis does not include a complete hardware
implementation of the upstream queue/SRAM logic that produces and transports
all descriptor inputs. As with next-SA descriptor generation, this interface
boundary is reported explicitly; the present area and power results must not
be interpreted as including every metadata-producer cost.

\subsection{Non-Retentive Power-Gating Safety}
\label{sec:nonretentive-safety}

Power-gating the SA is non-retentive: values held in the processing-element
registers are not assumed to survive an OFF interval. Consequently, an idle
interval alone is insufficient to authorize power gating. The controller may
enter the OFF state only when it has both (i) evidence that the SA cannot be
needed before a certified deadline and (ii) evidence that losing the current
SA-resident state cannot change subsequent execution.

The second requirement follows from the weight-stationary instruction
semantics. A \texttt{GEMM} instruction consumes weights already resident in
the SA and therefore cannot safely follow a non-retentive OFF interval. A
\texttt{GEMM\_PRELOAD} instruction reloads the required weights before their
use. Hence, for each of the two resident tile slots, the controller requires
the earliest future SA instruction to exist and to be
\texttt{GEMM\_PRELOAD}. It also requires both resident slots to be occupied;
the future-tile bound described in Sec.~\ref{sec:causal-predictor} limits the
certificate to the earliest point at which a resident tile could retire and
previously unseen work could enter.

Table~\ref{tab:nonretentive-safety} summarizes the hazards checked by the
implemented controller. These conditions are conjunctive: failure of any
entry condition suppresses a new OFF decision.

\begin{table}[t]
  \centering
  \caption{Safety conditions for a non-retentive SA OFF interval.}
  \label{tab:nonretentive-safety}
  \begin{tabular}{p{0.31\columnwidth}p{0.61\columnwidth}}
    \toprule
    Hazard & Implemented protection \\
    \midrule
    Useful SA work is already active
      & Require the SA to be inactive at every certificate stage and at
        entry commit. \\
    Committed work can reach the SA immediately
      & Require both the compute pipeline and execution-instruction queue to
        be empty. \\
    OFF destroys resident weights
      & Require the first future SA instruction in each resident tile to be
        \texttt{GEMM\_PRELOAD}, not \texttt{GEMM}. \\
    Unseen tile work arrives early
      & Require two resident tiles and cap the horizon using resident
        retirement and global work-drain information. \\
    Queue-distance evidence is ambiguous
      & Require exact memory-request tickets throughout certificate
        generation and at entry commit. \\
    Metadata changes after sampling
      & A write to either resident descriptor invalidates the certificate;
        a write during sleep forces wake-up. \\
    A saved deadline becomes stale
      & Freeze the committed deadline, decrement it monotonically, and never
        recompute or extend it while asleep. \\
    SA demand appears before the deadline
      & Wake immediately if the SA becomes active, either queue becomes
        nonempty, ticket exactness is lost, or the countdown reaches the
        one-cycle guard. \\
    \bottomrule
  \end{tabular}
\end{table}

Certificate construction is itself protected against stale observations.
The implementation uses three registered evaluation stages. A descriptor
write, observed SA activity, a nonempty compute or execution queue, or loss of
ticket exactness invalidates the certificate at the intervening stage
boundaries. The three elapsed evaluation cycles are deducted from the causal
lower bound. With the evaluated 166-cycle break-even threshold and one-cycle
wake guard, a fresh bound must therefore be at least 170 cycles; after the
three-stage deduction, entry commits only when at least 167 cycles remain.

Once entry commits, the controller stores the certified remaining interval in
a countdown register. While OFF, the controller cannot increase this value;
it either decrements the countdown or exits the OFF state. This frozen-deadline
rule is important because continuously recomputing a later deadline from new
observations could incorrectly prolong an interval that was justified by an
older machine state.

We validate the decision logic by replaying each accepted architectural trace
against both the reference controller and its clock-gated implementation. A
replay is accepted only when their emitted sleep episodes and summary records
match exactly. Independently, the replay testbench counts a safety violation
whenever the controller requests sleep while the trace reports the SA active,
and terminates with an error if this count is nonzero. Thus, the validation
establishes trace-level equivalence of the two controller implementations and
the absence of sleep--activity overlap in an accepted replay.

This validation has a deliberate scope. The controller is evaluated in
shadow replay and does not feed a physical OFF/ON transition back into
ONNXim's scheduling or timing. Therefore, unchanged workload cycle counts are
not, by themselves, a measurement of zero wake-up penalty in an integrated
power-gated NPU. Device-level transition delay and energy are characterized
separately, and an end-to-end implementation that stalls execution during a
late wake-up would be required to measure any resulting performance penalty.
